% 11_peripherals.tex -- APB bridge and bus, APB waveform, UART, SPI, I2C, GPIO / timer / watchdog.
% Build: tectonic 11_peripherals.tex
\documentclass[tikz,border=10pt,12pt]{standalone}
% soc_style.tex -- shared look for the SoC figures (same palette as the slide decks and the
% core register placement figure). Every figure does % soc_style.tex -- shared look for the SoC figures (same palette as the slide decks and the
% core register placement figure). Every figure does % soc_style.tex -- shared look for the SoC figures (same palette as the slide decks and the
% core register placement figure). Every figure does \input{soc_style}.
\usepackage[T1]{fontenc}
\usepackage[scaled=0.95]{helvet}
\renewcommand{\familydefault}{\sfdefault}
\usepackage{booktabs}
\usepackage{amsmath}
\usepackage{amssymb}
\usetikzlibrary{arrows.meta,calc,positioning,fit,shapes.geometric,shapes.misc}

% ---------------- palette ----------------
\definecolor{pcb}{HTML}{0B2E33}
\definecolor{pcbmid}{HTML}{1F5C61}
\definecolor{tint}{HTML}{E6F0EF}
\definecolor{copper}{HTML}{B5651D}
\definecolor{coppert}{HTML}{F6E8DC}
\definecolor{ink}{HTML}{1F2A2E}
\definecolor{muted}{HTML}{5B6B70}
\definecolor{hardf}{HTML}{DCE6F2}
\definecolor{harde}{HTML}{2F4A7A}
\definecolor{good}{HTML}{2E7D32}
\definecolor{bad}{HTML}{B3261E}
% one colour per kind of block: border (d) / fill (f)
\definecolor{cpud}{HTML}{2B5C9E} \definecolor{cpuf}{HTML}{D6E4F5}
\definecolor{gpud}{HTML}{2E7D32} \definecolor{gpuf}{HTML}{D5EDD0}
\definecolor{npud}{HTML}{0B2E33} \definecolor{npuf}{HTML}{BFE0DB}
\definecolor{memd}{HTML}{5E3A8C} \definecolor{memf}{HTML}{E4DAF3}
\definecolor{busd}{HTML}{B5651D} \definecolor{busf}{HTML}{F6E8DC}
\definecolor{prd}{HTML}{546E7A}  \definecolor{prf}{HTML}{E3EAED}
\definecolor{secd}{HTML}{B3261E} \definecolor{secf}{HTML}{F7D9D6}
\definecolor{pwrd}{HTML}{8A6800} \definecolor{pwrf}{HTML}{FFF0B8}
\definecolor{medd}{HTML}{8E2463} \definecolor{medf}{HTML}{F5D6E6}
\definecolor{skd}{HTML}{2E7D32}  \definecolor{skf}{HTML}{BFE3B4}   % skew buffers, as in the register figure

\tikzset{
  blk/.style 2 args={draw=#1, fill=#2, rounded corners=4pt, thick, align=center,
                     font=\scriptsize, inner sep=3pt},
  cpu/.style={blk={cpud}{cpuf}},  gpu/.style={blk={gpud}{gpuf}},
  npu/.style={blk={npud}{npuf}},  mem/.style={blk={memd}{memf}},
  bus/.style={blk={busd}{busf}},  per/.style={blk={prd}{prf}},
  sec/.style={blk={secd}{secf}},  pwr/.style={blk={pwrd}{pwrf}},
  med/.style={blk={medd}{medf}},
  comb/.style ={draw=copper, fill=coppert, rounded corners=2pt, thick, align=center,
                font=\scriptsize, inner sep=2pt},                       % combinational logic
  hard/.style ={draw=harde, fill=hardf, thick, dashed, align=center,
                font=\scriptsize, inner sep=2pt},                       % hard macro / off-chip
  sram/.style ={draw=memd, fill=memf, thick, align=center, font=\scriptsize,
                inner sep=2pt},                                         % SRAM array, not flip-flops
  ff/.style   ={draw=#1, fill=#1!30, line width=0.35pt},                % one flip-flop square
  ffgone/.style={draw=muted, fill=white, densely dashed, line width=0.35pt},
  panel/.style={draw=muted, rounded corners=6pt, line width=0.9pt},
  ptitle/.style={anchor=north west, font=\normalsize\bfseries, text=ink, inner sep=0pt},
  note/.style ={font=\scriptsize, text=ink, align=left, inner sep=0pt},
  lab/.style  ={font=\tiny, text=muted, inner sep=1pt},
  axi/.style  ={{Stealth[length=2.6mm]}-{Stealth[length=2.6mm]}, line width=2.2pt, copper},
  axi1/.style ={-{Stealth[length=2.6mm]}, line width=2.2pt, copper},
  apb/.style  ={{Stealth[length=1.8mm]}-{Stealth[length=1.8mm]}, line width=1pt, copper},
  apb1/.style ={-{Stealth[length=1.8mm]}, line width=1pt, copper},
  dat/.style  ={-{Stealth[length=1.8mm]}, thick, pcbmid},
  ctl/.style  ={-{Stealth[length=1.8mm]}, densely dashed, thick, muted},
  irq/.style  ={-{Stealth[length=2mm]}, densely dotted, line width=1.3pt, bad},
  num/.style  ={circle, fill=copper, text=white, font=\scriptsize\bfseries, inner sep=1pt,
                minimum size=0.46cm},
  ours/.style ={draw=bad, line width=1.2pt, densely dashed, rounded corners=4pt},
}

% ---------------- helpers ----------------
\newcommand{\rn}[1]{\texttt{#1}}
% \ffrow{x}{y}{n}{cell}{style}: n flip-flops in a row, bottom-left corner at (x,y)
\newcommand{\ffrow}[5]{\foreach \q in {1,...,#3}{\draw[#5] ({#1+(\q-1)*#4},{#2}) rectangle ++(#4,#4);}}
% \ffcol{x}{y}{n}{cell}{style}: n flip-flops in a column, top-left corner at (x,y)
\newcommand{\ffcol}[5]{\foreach \q in {1,...,#3}{\draw[#5] ({#1},{#2-\q*#4}) rectangle ++(#4,#4);}}
% \ffgrid{x}{ytop}{n}{per row}{cell}{style}: n flip-flops, filled row by row from the top-left
\newcommand{\ffgrid}[6]{\pgfmathtruncatemacro{\gl}{#3-1}%
  \foreach \b in {0,...,\gl}{\pgfmathtruncatemacro{\gc}{mod(\b,#4)}\pgfmathtruncatemacro{\gr}{div(\b,#4)}%
    \draw[#6] ({#1+\gc*#5},{#2-(\gr+1)*#5}) rectangle ++(#5,#5);}}

% every figure: same page size, title, footer
\newcommand{\canvas}{\useasboundingbox (-0.3,-0.2) rectangle (38.6,30.4);}
\newcommand{\figtitle}[3]{%
  \node[num, minimum size=0.9cm, font=\Large\bfseries] at (0.45,29.55) {#1};
  \node[anchor=west, font=\LARGE\bfseries, text=pcb] at (1.15,29.55) {#2};
  \node[anchor=west, font=\small, text=ink] at (0,28.7) {#3};}
\newcommand{\figfoot}[1]{\node[anchor=west, font=\scriptsize, text=muted] at (0,0.25) {#1};}
% \panel{x0}{y0}{x1}{y1}{title}
\newcommand{\panel}[5]{\draw[panel] (#1,#2) rectangle (#3,#4);
  \node[ptitle] at ({#1+0.3},{#4-0.25}) {#5};}

% ---------------- timing diagrams: 1 cycle = \cw cm, low at y, high at y+0.45 ----------------
\newcommand{\cw}{1.0}
\newcommand{\wclk}[3]{% x0 y ncycles
  \foreach \k in {1,...,#3}{\draw[pcb, thick] ({#1+(\k-1)*\cw},#2) -- ++(0,0.45) -- ++(0.5*\cw,0)
     -- ++(0,-0.45) -- ++(0.5*\cw,0);}}
\newcommand{\wsig}[3]{% x0 y {level per cycle, e.g. 0,1,1,0}
  \gdef\wprev{x}%
  \foreach \v [count=\k from 0] in {#3}{%
    \draw[pcb, thick] ({#1+\k*\cw},{#2+0.45*\v}) -- ++(\cw,0);
    \if x\wprev\else\ifnum\wprev=\v\else\draw[pcb, thick] ({#1+\k*\cw},#2) -- ++(0,0.45);\fi\fi
    \global\let\wprev\v}}
\newcommand{\wbus}[5]{% x0 y c0 c1 text: one value held from cycle c0 to cycle c1
  \draw[pcbmid, thick, fill=white] ({#1+(#3)*\cw},{#2+0.225}) -- ({#1+(#3)*\cw+0.08},{#2+0.45})
    -- ({#1+(#4)*\cw-0.08},{#2+0.45}) -- ({#1+(#4)*\cw},{#2+0.225}) -- ({#1+(#4)*\cw-0.08},#2)
    -- ({#1+(#3)*\cw+0.08},#2) -- cycle;
  \node[font=\tiny\ttfamily] at ({#1+((#3)+(#4))*\cw/2},{#2+0.225}) {#5};}
\newcommand{\widle}[4]{% x0 y c0 c1: bus idle (no value)
  \draw[muted, thick] ({#1+(#3)*\cw},{#2+0.225}) -- ({#1+(#4)*\cw},{#2+0.225});}
\newcommand{\wname}[3]{\node[anchor=east, font=\scriptsize\ttfamily] at ({#1-0.15},{#2+0.225}) {#3};}
.
\usepackage[T1]{fontenc}
\usepackage[scaled=0.95]{helvet}
\renewcommand{\familydefault}{\sfdefault}
\usepackage{booktabs}
\usepackage{amsmath}
\usepackage{amssymb}
\usetikzlibrary{arrows.meta,calc,positioning,fit,shapes.geometric,shapes.misc}

% ---------------- palette ----------------
\definecolor{pcb}{HTML}{0B2E33}
\definecolor{pcbmid}{HTML}{1F5C61}
\definecolor{tint}{HTML}{E6F0EF}
\definecolor{copper}{HTML}{B5651D}
\definecolor{coppert}{HTML}{F6E8DC}
\definecolor{ink}{HTML}{1F2A2E}
\definecolor{muted}{HTML}{5B6B70}
\definecolor{hardf}{HTML}{DCE6F2}
\definecolor{harde}{HTML}{2F4A7A}
\definecolor{good}{HTML}{2E7D32}
\definecolor{bad}{HTML}{B3261E}
% one colour per kind of block: border (d) / fill (f)
\definecolor{cpud}{HTML}{2B5C9E} \definecolor{cpuf}{HTML}{D6E4F5}
\definecolor{gpud}{HTML}{2E7D32} \definecolor{gpuf}{HTML}{D5EDD0}
\definecolor{npud}{HTML}{0B2E33} \definecolor{npuf}{HTML}{BFE0DB}
\definecolor{memd}{HTML}{5E3A8C} \definecolor{memf}{HTML}{E4DAF3}
\definecolor{busd}{HTML}{B5651D} \definecolor{busf}{HTML}{F6E8DC}
\definecolor{prd}{HTML}{546E7A}  \definecolor{prf}{HTML}{E3EAED}
\definecolor{secd}{HTML}{B3261E} \definecolor{secf}{HTML}{F7D9D6}
\definecolor{pwrd}{HTML}{8A6800} \definecolor{pwrf}{HTML}{FFF0B8}
\definecolor{medd}{HTML}{8E2463} \definecolor{medf}{HTML}{F5D6E6}
\definecolor{skd}{HTML}{2E7D32}  \definecolor{skf}{HTML}{BFE3B4}   % skew buffers, as in the register figure

\tikzset{
  blk/.style 2 args={draw=#1, fill=#2, rounded corners=4pt, thick, align=center,
                     font=\scriptsize, inner sep=3pt},
  cpu/.style={blk={cpud}{cpuf}},  gpu/.style={blk={gpud}{gpuf}},
  npu/.style={blk={npud}{npuf}},  mem/.style={blk={memd}{memf}},
  bus/.style={blk={busd}{busf}},  per/.style={blk={prd}{prf}},
  sec/.style={blk={secd}{secf}},  pwr/.style={blk={pwrd}{pwrf}},
  med/.style={blk={medd}{medf}},
  comb/.style ={draw=copper, fill=coppert, rounded corners=2pt, thick, align=center,
                font=\scriptsize, inner sep=2pt},                       % combinational logic
  hard/.style ={draw=harde, fill=hardf, thick, dashed, align=center,
                font=\scriptsize, inner sep=2pt},                       % hard macro / off-chip
  sram/.style ={draw=memd, fill=memf, thick, align=center, font=\scriptsize,
                inner sep=2pt},                                         % SRAM array, not flip-flops
  ff/.style   ={draw=#1, fill=#1!30, line width=0.35pt},                % one flip-flop square
  ffgone/.style={draw=muted, fill=white, densely dashed, line width=0.35pt},
  panel/.style={draw=muted, rounded corners=6pt, line width=0.9pt},
  ptitle/.style={anchor=north west, font=\normalsize\bfseries, text=ink, inner sep=0pt},
  note/.style ={font=\scriptsize, text=ink, align=left, inner sep=0pt},
  lab/.style  ={font=\tiny, text=muted, inner sep=1pt},
  axi/.style  ={{Stealth[length=2.6mm]}-{Stealth[length=2.6mm]}, line width=2.2pt, copper},
  axi1/.style ={-{Stealth[length=2.6mm]}, line width=2.2pt, copper},
  apb/.style  ={{Stealth[length=1.8mm]}-{Stealth[length=1.8mm]}, line width=1pt, copper},
  apb1/.style ={-{Stealth[length=1.8mm]}, line width=1pt, copper},
  dat/.style  ={-{Stealth[length=1.8mm]}, thick, pcbmid},
  ctl/.style  ={-{Stealth[length=1.8mm]}, densely dashed, thick, muted},
  irq/.style  ={-{Stealth[length=2mm]}, densely dotted, line width=1.3pt, bad},
  num/.style  ={circle, fill=copper, text=white, font=\scriptsize\bfseries, inner sep=1pt,
                minimum size=0.46cm},
  ours/.style ={draw=bad, line width=1.2pt, densely dashed, rounded corners=4pt},
}

% ---------------- helpers ----------------
\newcommand{\rn}[1]{\texttt{#1}}
% \ffrow{x}{y}{n}{cell}{style}: n flip-flops in a row, bottom-left corner at (x,y)
\newcommand{\ffrow}[5]{\foreach \q in {1,...,#3}{\draw[#5] ({#1+(\q-1)*#4},{#2}) rectangle ++(#4,#4);}}
% \ffcol{x}{y}{n}{cell}{style}: n flip-flops in a column, top-left corner at (x,y)
\newcommand{\ffcol}[5]{\foreach \q in {1,...,#3}{\draw[#5] ({#1},{#2-\q*#4}) rectangle ++(#4,#4);}}
% \ffgrid{x}{ytop}{n}{per row}{cell}{style}: n flip-flops, filled row by row from the top-left
\newcommand{\ffgrid}[6]{\pgfmathtruncatemacro{\gl}{#3-1}%
  \foreach \b in {0,...,\gl}{\pgfmathtruncatemacro{\gc}{mod(\b,#4)}\pgfmathtruncatemacro{\gr}{div(\b,#4)}%
    \draw[#6] ({#1+\gc*#5},{#2-(\gr+1)*#5}) rectangle ++(#5,#5);}}

% every figure: same page size, title, footer
\newcommand{\canvas}{\useasboundingbox (-0.3,-0.2) rectangle (38.6,30.4);}
\newcommand{\figtitle}[3]{%
  \node[num, minimum size=0.9cm, font=\Large\bfseries] at (0.45,29.55) {#1};
  \node[anchor=west, font=\LARGE\bfseries, text=pcb] at (1.15,29.55) {#2};
  \node[anchor=west, font=\small, text=ink] at (0,28.7) {#3};}
\newcommand{\figfoot}[1]{\node[anchor=west, font=\scriptsize, text=muted] at (0,0.25) {#1};}
% \panel{x0}{y0}{x1}{y1}{title}
\newcommand{\panel}[5]{\draw[panel] (#1,#2) rectangle (#3,#4);
  \node[ptitle] at ({#1+0.3},{#4-0.25}) {#5};}

% ---------------- timing diagrams: 1 cycle = \cw cm, low at y, high at y+0.45 ----------------
\newcommand{\cw}{1.0}
\newcommand{\wclk}[3]{% x0 y ncycles
  \foreach \k in {1,...,#3}{\draw[pcb, thick] ({#1+(\k-1)*\cw},#2) -- ++(0,0.45) -- ++(0.5*\cw,0)
     -- ++(0,-0.45) -- ++(0.5*\cw,0);}}
\newcommand{\wsig}[3]{% x0 y {level per cycle, e.g. 0,1,1,0}
  \gdef\wprev{x}%
  \foreach \v [count=\k from 0] in {#3}{%
    \draw[pcb, thick] ({#1+\k*\cw},{#2+0.45*\v}) -- ++(\cw,0);
    \if x\wprev\else\ifnum\wprev=\v\else\draw[pcb, thick] ({#1+\k*\cw},#2) -- ++(0,0.45);\fi\fi
    \global\let\wprev\v}}
\newcommand{\wbus}[5]{% x0 y c0 c1 text: one value held from cycle c0 to cycle c1
  \draw[pcbmid, thick, fill=white] ({#1+(#3)*\cw},{#2+0.225}) -- ({#1+(#3)*\cw+0.08},{#2+0.45})
    -- ({#1+(#4)*\cw-0.08},{#2+0.45}) -- ({#1+(#4)*\cw},{#2+0.225}) -- ({#1+(#4)*\cw-0.08},#2)
    -- ({#1+(#3)*\cw+0.08},#2) -- cycle;
  \node[font=\tiny\ttfamily] at ({#1+((#3)+(#4))*\cw/2},{#2+0.225}) {#5};}
\newcommand{\widle}[4]{% x0 y c0 c1: bus idle (no value)
  \draw[muted, thick] ({#1+(#3)*\cw},{#2+0.225}) -- ({#1+(#4)*\cw},{#2+0.225});}
\newcommand{\wname}[3]{\node[anchor=east, font=\scriptsize\ttfamily] at ({#1-0.15},{#2+0.225}) {#3};}
.
\usepackage[T1]{fontenc}
\usepackage[scaled=0.95]{helvet}
\renewcommand{\familydefault}{\sfdefault}
\usepackage{booktabs}
\usepackage{amsmath}
\usepackage{amssymb}
\usetikzlibrary{arrows.meta,calc,positioning,fit,shapes.geometric,shapes.misc}

% ---------------- palette ----------------
\definecolor{pcb}{HTML}{0B2E33}
\definecolor{pcbmid}{HTML}{1F5C61}
\definecolor{tint}{HTML}{E6F0EF}
\definecolor{copper}{HTML}{B5651D}
\definecolor{coppert}{HTML}{F6E8DC}
\definecolor{ink}{HTML}{1F2A2E}
\definecolor{muted}{HTML}{5B6B70}
\definecolor{hardf}{HTML}{DCE6F2}
\definecolor{harde}{HTML}{2F4A7A}
\definecolor{good}{HTML}{2E7D32}
\definecolor{bad}{HTML}{B3261E}
% one colour per kind of block: border (d) / fill (f)
\definecolor{cpud}{HTML}{2B5C9E} \definecolor{cpuf}{HTML}{D6E4F5}
\definecolor{gpud}{HTML}{2E7D32} \definecolor{gpuf}{HTML}{D5EDD0}
\definecolor{npud}{HTML}{0B2E33} \definecolor{npuf}{HTML}{BFE0DB}
\definecolor{memd}{HTML}{5E3A8C} \definecolor{memf}{HTML}{E4DAF3}
\definecolor{busd}{HTML}{B5651D} \definecolor{busf}{HTML}{F6E8DC}
\definecolor{prd}{HTML}{546E7A}  \definecolor{prf}{HTML}{E3EAED}
\definecolor{secd}{HTML}{B3261E} \definecolor{secf}{HTML}{F7D9D6}
\definecolor{pwrd}{HTML}{8A6800} \definecolor{pwrf}{HTML}{FFF0B8}
\definecolor{medd}{HTML}{8E2463} \definecolor{medf}{HTML}{F5D6E6}
\definecolor{skd}{HTML}{2E7D32}  \definecolor{skf}{HTML}{BFE3B4}   % skew buffers, as in the register figure

\tikzset{
  blk/.style 2 args={draw=#1, fill=#2, rounded corners=4pt, thick, align=center,
                     font=\scriptsize, inner sep=3pt},
  cpu/.style={blk={cpud}{cpuf}},  gpu/.style={blk={gpud}{gpuf}},
  npu/.style={blk={npud}{npuf}},  mem/.style={blk={memd}{memf}},
  bus/.style={blk={busd}{busf}},  per/.style={blk={prd}{prf}},
  sec/.style={blk={secd}{secf}},  pwr/.style={blk={pwrd}{pwrf}},
  med/.style={blk={medd}{medf}},
  comb/.style ={draw=copper, fill=coppert, rounded corners=2pt, thick, align=center,
                font=\scriptsize, inner sep=2pt},                       % combinational logic
  hard/.style ={draw=harde, fill=hardf, thick, dashed, align=center,
                font=\scriptsize, inner sep=2pt},                       % hard macro / off-chip
  sram/.style ={draw=memd, fill=memf, thick, align=center, font=\scriptsize,
                inner sep=2pt},                                         % SRAM array, not flip-flops
  ff/.style   ={draw=#1, fill=#1!30, line width=0.35pt},                % one flip-flop square
  ffgone/.style={draw=muted, fill=white, densely dashed, line width=0.35pt},
  panel/.style={draw=muted, rounded corners=6pt, line width=0.9pt},
  ptitle/.style={anchor=north west, font=\normalsize\bfseries, text=ink, inner sep=0pt},
  note/.style ={font=\scriptsize, text=ink, align=left, inner sep=0pt},
  lab/.style  ={font=\tiny, text=muted, inner sep=1pt},
  axi/.style  ={{Stealth[length=2.6mm]}-{Stealth[length=2.6mm]}, line width=2.2pt, copper},
  axi1/.style ={-{Stealth[length=2.6mm]}, line width=2.2pt, copper},
  apb/.style  ={{Stealth[length=1.8mm]}-{Stealth[length=1.8mm]}, line width=1pt, copper},
  apb1/.style ={-{Stealth[length=1.8mm]}, line width=1pt, copper},
  dat/.style  ={-{Stealth[length=1.8mm]}, thick, pcbmid},
  ctl/.style  ={-{Stealth[length=1.8mm]}, densely dashed, thick, muted},
  irq/.style  ={-{Stealth[length=2mm]}, densely dotted, line width=1.3pt, bad},
  num/.style  ={circle, fill=copper, text=white, font=\scriptsize\bfseries, inner sep=1pt,
                minimum size=0.46cm},
  ours/.style ={draw=bad, line width=1.2pt, densely dashed, rounded corners=4pt},
}

% ---------------- helpers ----------------
\newcommand{\rn}[1]{\texttt{#1}}
% \ffrow{x}{y}{n}{cell}{style}: n flip-flops in a row, bottom-left corner at (x,y)
\newcommand{\ffrow}[5]{\foreach \q in {1,...,#3}{\draw[#5] ({#1+(\q-1)*#4},{#2}) rectangle ++(#4,#4);}}
% \ffcol{x}{y}{n}{cell}{style}: n flip-flops in a column, top-left corner at (x,y)
\newcommand{\ffcol}[5]{\foreach \q in {1,...,#3}{\draw[#5] ({#1},{#2-\q*#4}) rectangle ++(#4,#4);}}
% \ffgrid{x}{ytop}{n}{per row}{cell}{style}: n flip-flops, filled row by row from the top-left
\newcommand{\ffgrid}[6]{\pgfmathtruncatemacro{\gl}{#3-1}%
  \foreach \b in {0,...,\gl}{\pgfmathtruncatemacro{\gc}{mod(\b,#4)}\pgfmathtruncatemacro{\gr}{div(\b,#4)}%
    \draw[#6] ({#1+\gc*#5},{#2-(\gr+1)*#5}) rectangle ++(#5,#5);}}

% every figure: same page size, title, footer
\newcommand{\canvas}{\useasboundingbox (-0.3,-0.2) rectangle (38.6,30.4);}
\newcommand{\figtitle}[3]{%
  \node[num, minimum size=0.9cm, font=\Large\bfseries] at (0.45,29.55) {#1};
  \node[anchor=west, font=\LARGE\bfseries, text=pcb] at (1.15,29.55) {#2};
  \node[anchor=west, font=\small, text=ink] at (0,28.7) {#3};}
\newcommand{\figfoot}[1]{\node[anchor=west, font=\scriptsize, text=muted] at (0,0.25) {#1};}
% \panel{x0}{y0}{x1}{y1}{title}
\newcommand{\panel}[5]{\draw[panel] (#1,#2) rectangle (#3,#4);
  \node[ptitle] at ({#1+0.3},{#4-0.25}) {#5};}

% ---------------- timing diagrams: 1 cycle = \cw cm, low at y, high at y+0.45 ----------------
\newcommand{\cw}{1.0}
\newcommand{\wclk}[3]{% x0 y ncycles
  \foreach \k in {1,...,#3}{\draw[pcb, thick] ({#1+(\k-1)*\cw},#2) -- ++(0,0.45) -- ++(0.5*\cw,0)
     -- ++(0,-0.45) -- ++(0.5*\cw,0);}}
\newcommand{\wsig}[3]{% x0 y {level per cycle, e.g. 0,1,1,0}
  \gdef\wprev{x}%
  \foreach \v [count=\k from 0] in {#3}{%
    \draw[pcb, thick] ({#1+\k*\cw},{#2+0.45*\v}) -- ++(\cw,0);
    \if x\wprev\else\ifnum\wprev=\v\else\draw[pcb, thick] ({#1+\k*\cw},#2) -- ++(0,0.45);\fi\fi
    \global\let\wprev\v}}
\newcommand{\wbus}[5]{% x0 y c0 c1 text: one value held from cycle c0 to cycle c1
  \draw[pcbmid, thick, fill=white] ({#1+(#3)*\cw},{#2+0.225}) -- ({#1+(#3)*\cw+0.08},{#2+0.45})
    -- ({#1+(#4)*\cw-0.08},{#2+0.45}) -- ({#1+(#4)*\cw},{#2+0.225}) -- ({#1+(#4)*\cw-0.08},#2)
    -- ({#1+(#3)*\cw+0.08},#2) -- cycle;
  \node[font=\tiny\ttfamily] at ({#1+((#3)+(#4))*\cw/2},{#2+0.225}) {#5};}
\newcommand{\widle}[4]{% x0 y c0 c1: bus idle (no value)
  \draw[muted, thick] ({#1+(#3)*\cw},{#2+0.225}) -- ({#1+(#4)*\cw},{#2+0.225});}
\newcommand{\wname}[3]{\node[anchor=east, font=\scriptsize\ttfamily] at ({#1-0.15},{#2+0.225}) {#3};}

\newcommand{\tick}[2]{\node[circle, fill=good, text=white, font=\tiny\bfseries, inner sep=0.5pt] at (#1,#2) {\checkmark};}
\begin{document}
\begin{tikzpicture}[line join=round]
\canvas
\figtitle{11}{Peripherals: slow devices on the APB bus}
  {UART, SPI, I2C, GPIO and timers are small blocks: a few registers plus a shift register or a counter.
   The CPU reaches them through an AXI-to-APB bridge, and each one has an IRQ wire to the GIC.}

% =====================================================================
% (A) bridge + bus
% =====================================================================
\panel{0}{15.0}{14.0}{28.1}{AXI $\rightarrow$ APB bridge and the peripheral bus}
\node[bus, minimum width=13.3cm, minimum height=0.75cm, anchor=north west] at (0.35,27.3) {AXI from the NoC (fast, 128 bit)};
\node[bus, minimum width=13.3cm, minimum height=1.35cm, anchor=north west, text width=13.0cm] at (0.35,26.3)
  {\textbf{AXI $\rightarrow$ APB bridge}: a small FSM; turns one AXI request into one APB transfer
   and decodes PADDR $\rightarrow$ one PSEL per peripheral};
\draw[copper, line width=1.4pt] (0.8,24.0) -- (13.4,24.0);
\node[lab, text=busd, anchor=south west, font=\tiny\bfseries] at (0.8,24.02) {shared: PADDR, PWDATA, PWRITE, PENABLE (32 bit, e.g. 100 MHz)};
\foreach \nm/\ad [count=\k from 0] in {UART/0x1000\_0000, SPI/0x1000\_1000, I2C/0x1000\_2000, GPIO/0x1000\_3000, timer/0x1000\_4000, watchdog/0x1000\_5000}{
  \pgfmathsetmacro{\x}{0.35+\k*2.25}
  \node[per, fill=white, minimum width=2.05cm, minimum height=1.9cm, text width=1.9cm, anchor=north west, font=\tiny]
    (p\k) at (\x,23.0) {\textbf{\nm}\\[2pt]\ad};
  \draw[apb1] ({\x+1.0},24.0) -- ({\x+1.0},23.0);
  \node[lab, anchor=west] at ({\x+1.05},23.5) {PSEL\k};
  \draw[irq] ({\x+1.0},21.1) -- ({\x+1.0},20.45);}
\draw[bad, densely dotted, line width=1.3pt] (1.35,20.45) -- (12.6,20.45);
\node[font=\tiny\bfseries, text=bad, anchor=north] at (7.0,20.4) {6 IRQ wires $\rightarrow$ GIC};
\node[note, anchor=north west, text width=13.3cm] at (0.35,19.6) {%
  Example map: 4 KB per peripheral. Inside each one, the register offsets work exactly like your \rn{mm\_regs}:
  \rn{0x00} DATA, \rn{0x04} STATUS, \rn{0x08} CTRL, \ldots\\[3pt]
  \textbf{Why a separate slow bus?} Peripherals need only a few transfers per second. APB has no bursts and no IDs, so each
  peripheral's bus interface is a handful of gates, and the fast AXI fabric is not slowed down by them.};

% =====================================================================
% (B) APB waveform
% =====================================================================
\panel{14.3}{21.2}{38.3}{28.1}{APB: a write, then a read with one wait cycle}
\renewcommand{\cw}{1.3}
\foreach \c in {0,...,9}{\node[lab] at ({17.2+(\c+0.5)*\cw},27.35) {\c};}
\wname{17.2}{26.65}{PCLK}     \wclk{17.2}{26.65}{10}
\wname{17.2}{26.0}{PSEL}      \wsig{17.2}{26.0}{0,1,1,0,1,1,1,0,0,0}
\wname{17.2}{25.35}{PENABLE}  \wsig{17.2}{25.35}{0,0,1,0,0,1,1,0,0,0}
\wname{17.2}{24.7}{PWRITE}    \wsig{17.2}{24.7}{0,1,1,0,0,0,0,0,0,0}
\wname{17.2}{24.05}{PADDR}
\widle{17.2}{24.05}{0}{1} \wbus{17.2}{24.05}{1}{3}{0x00 DATA} \widle{17.2}{24.05}{3}{4} \wbus{17.2}{24.05}{4}{7}{0x04 STATUS} \widle{17.2}{24.05}{7}{10}
\wname{17.2}{23.4}{PWDATA}
\widle{17.2}{23.4}{0}{1} \wbus{17.2}{23.4}{1}{3}{0x42} \widle{17.2}{23.4}{3}{10}
\wname{17.2}{22.75}{PREADY}   \wsig{17.2}{22.75}{0,0,1,0,0,0,1,0,0,0}
\wname{17.2}{22.1}{PRDATA}
\widle{17.2}{22.1}{0}{6} \wbus{17.2}{22.1}{6}{7}{0x01} \widle{17.2}{22.1}{7}{10}
\foreach \c in {3,7}{\draw[good, thick, densely dashed] ({17.2+\c*\cw},22.0) -- ({17.2+\c*\cw},27.15); \tick{{17.2+\c*\cw}}{21.75}}
\node[note, font=\tiny, anchor=north west, text width=6.5cm] at (31.6,27.3) {%
  \textbf{Setup phase}: PSEL = 1, PENABLE = 0, address and data on the bus.\\[3pt]
  \textbf{Access phase}: PENABLE = 1. The transfer ends on the first edge with PREADY = 1.\\[3pt]
  Cycle 5: the peripheral is not ready (PREADY = 0), so the access phase is stretched by one cycle.\\[3pt]
  Write 0x42 to UART DATA = send the byte 0x42. Read STATUS = 0x01 (TX busy).\\[3pt]
  Minimum 2 cycles per transfer; no overlapping.};

% =====================================================================
% (C) UART
% =====================================================================
\panel{14.3}{15.0}{38.3}{20.9}{UART: one byte as a frame of 10 bits (your \rn{uart\_tx.sv} / \rn{uart\_rx.sv})}
\node[anchor=north west, font=\tiny\bfseries] at (14.6,20.15) {registers};
\foreach \nm/\n [count=\k from 0] in {DATA/8, STATUS/2, CTRL/2, BAUD\_DIV/16}{
  \node[font=\tiny\ttfamily, anchor=west] at (14.6,{19.55-\k*0.5}) {\nm};
  \ffrow{16.2}{19.43-\k*0.5}{\n}{0.2}{ff=prd}}
\node[lab, anchor=west] at (19.5,18.06) {= 434};
\node[anchor=north west, font=\tiny\bfseries] at (14.6,17.45) {TX shift register (10 bits): stop, d7 \ldots d0, start};
\ffrow{14.65}{16.6}{10}{0.36}{ff=prd}
\foreach \t [count=\k from 0] in {1,0,1,0,0,0,0,1,0,0}{\node[font=\tiny] at ({14.83+\k*0.36},16.78) {\t};}
\draw[dat] (18.25,16.78) -- (18.8,16.78); \node[lab, text=ink, anchor=west] at (18.8,16.78) {TX pin};
\node[note, font=\tiny, anchor=north west, text width=5.3cm] at (14.6,16.35) {%
  A counter divides 50 MHz by 434 (= 50,000,000 / 115,200, rounded): every 434 clocks the register shifts one bit out.};
\renewcommand{\cw}{1.2}
\foreach \nm [count=\k from 0] in {idle, start, d0, d1, d2, d3, d4, d5, d6, d7, stop, idle}{\node[lab] at ({23.6+(\k+0.5)*\cw},20.0) {\nm};}
\wname{23.6}{19.2}{TX}  \wsig{23.6}{19.2}{1,0,0,1,0,0,0,0,1,0,1,1}
\foreach \b [count=\k from 2] in {0,1,0,0,0,0,1,0}{\node[font=\tiny\bfseries, text=cpud] at ({23.6+(\k+0.5)*\cw},18.85) {\b};}
\draw[{Stealth[length=1.5mm]}-{Stealth[length=1.5mm]}, bad] ({23.6+1*\cw},18.45) -- ({23.6+2*\cw},18.45);
\node[font=\tiny, text=bad, anchor=north] at ({23.6+1.5*\cw},18.4) {8.68 $\mu$s};
\node[note, font=\tiny, anchor=north west, text width=14.0cm] at (23.6,17.85) {%
  Byte 0x42 = 0100\,0010, sent \textbf{LSB first}: 0,1,0,0,0,0,1,0. A start bit (0) wakes the receiver; a stop bit (1)
  returns the line to idle. 10 bits $\times$ 8.68 $\mu$s = 86.8 $\mu$s per byte $\Rightarrow$ 11,520 bytes/s at 115,200 baud.\\[3pt]
  The receiver syncs the pin with 2 flip-flops, waits for the falling edge, moves to the middle of the start bit
  (217 clocks), then samples every 434 clocks, always in the middle of a bit. This is exactly what your \rn{uart\_rx.sv} does.};

% =====================================================================
% (D) SPI
% =====================================================================
\panel{0}{0.7}{12.6}{14.7}{SPI: 4 wires, two shift registers in a ring}
\node[anchor=north west, font=\scriptsize\bfseries] at (0.3,13.95) {controller (SoC)};
\node[anchor=north west, font=\scriptsize\bfseries] at (7.3,13.95) {device (e.g. flash)};
\ffrow{0.4}{12.75}{8}{0.45}{ff=prd} \ffrow{7.4}{12.75}{8}{0.45}{ff=prd}
\draw[dat] (4.0,13.0) -- (7.4,13.0); \node[lab, text=ink, anchor=south] at (5.7,13.0) {MOSI};
\draw[dat] (11.0,12.97) -- (11.3,12.97) -- (11.3,12.25) -- (0.1,12.25) -- (0.1,12.97) -- (0.4,12.97);
\node[lab, text=ink, anchor=north] at (5.7,12.22) {MISO};
\node[note, font=\scriptsize, anchor=north west, text width=12.0cm] at (0.3,11.75) {%
  Every SCLK pulse both sides shift by one bit: after 8 pulses they have \textbf{swapped} bytes. CS\_n = 0 selects the device.};
\renewcommand{\cw}{1.0}
\wname{1.6}{10.0}{CS\_n}  \wsig{1.6}{10.0}{1,0,0,0,0,0,0,0,0,1}
\wname{1.6}{9.25}{SCLK}
\wsig{1.6}{9.25}{0,0,0,0,0,0,0,0,0,0}
\foreach \k in {1,...,8}{\draw[pcb, thick, fill=white] ({1.6+\k*\cw+0.25},9.25) -- ++(0,0.45) -- ++(0.5,0) -- ++(0,-0.45);}
\wname{1.6}{8.5}{MOSI}
\widle{1.6}{8.5}{0}{1}
\foreach \b [count=\k from 1] in {1,0,1,0,0,1,0,1}{\wbus{1.6}{8.5}{\k}{\k+1}{\b}}
\widle{1.6}{8.5}{9}{10}
\wname{1.6}{7.75}{MISO}
\widle{1.6}{7.75}{0}{1}
\foreach \b [count=\k from 1] in {0,0,1,1,1,1,0,0}{\wbus{1.6}{7.75}{\k}{\k+1}{\b}}
\widle{1.6}{7.75}{9}{10}
\node[note, font=\scriptsize, anchor=north west, text width=12.0cm] at (0.3,7.25) {%
  Mode 0: SCLK idles low, both sides sample on the rising edge, MSB first. MOSI sends 0xA5, MISO returns 0x3C.\\[3pt]
  Fast (1--50 MHz), full duplex, but one CS\_n wire per device. Used for flash chips, displays, ADCs and sensors.};

% =====================================================================
% (E) I2C
% =====================================================================
\panel{12.9}{0.7}{25.4}{14.7}{I2C: 2 shared wires, addressed devices}
\draw[pcb, thick] (13.6,13.6) -- (24.8,13.6); \node[lab, text=ink, anchor=south west] at (13.6,13.62) {VDD};
\foreach \x/\nm in {14.4/SDA, 15.4/SCL}{
  \draw[pcb, thick] (\x,13.6) -- (\x,13.3);
  \draw[pcb, thick] ({\x-0.12},13.3) rectangle ({\x+0.12},12.6);
  \draw[pcb, thick] (\x,12.6) -- (\x,12.2);
  \node[lab, text=ink, anchor=east] at ({\x-0.15},12.95) {R};}
\draw[pcb, thick] (14.4,12.2) -- (24.8,12.2); \node[lab, text=ink, anchor=west] at (24.8,12.2) {SDA};
\draw[pcb, thick] (15.4,11.7) -- (24.8,11.7); \node[lab, text=ink, anchor=west] at (24.8,11.7) {SCL};
\draw[pcb, thick] (15.4,12.2) -- (15.4,11.7);
\foreach \x/\nm in {17.0/{SoC}, 19.6/{sensor 0x1D}, 22.2/{EEPROM 0x50}}{
  \node[per, fill=white, minimum width=2.2cm, minimum height=0.75cm, font=\tiny, anchor=north] at (\x,11.25) {\nm};
  \draw[pcb] ({\x-0.4},12.2) -- ({\x-0.4},11.25); \draw[pcb] ({\x+0.4},11.7) -- ({\x+0.4},11.25);}
\node[note, font=\scriptsize, anchor=north west, text width=12.0cm] at (13.2,10.25) {%
  \textbf{Open drain}: a device can only pull a wire to 0; the resistor pulls it back to 1. So the wire is the AND of all
  devices and nobody can short it.};
\renewcommand{\cw}{0.82}
\foreach \nm [count=\k from 0] in {S, A6, A5, A4, A3, A2, A1, A0, R/W, ACK, P}{\node[lab] at ({14.6+(\k+0.5)*\cw},8.75) {\nm};}
\renewcommand{\cw}{0.41}
\wname{14.6}{8.0}{SCL} \wsig{14.6}{8.0}{1,1,0,1,0,1,0,1,0,1,0,1,0,1,0,1,0,1,0,1,1,1}
\wname{14.6}{7.2}{SDA} \wsig{14.6}{7.2}{1,0,0,0,0,0,1,1,1,1,1,1,0,0,1,1,0,0,0,0,0,1}
\renewcommand{\cw}{0.82}
\foreach \b [count=\k from 1] in {0,0,1,1,1,0,1,0,0}{\node[font=\tiny\bfseries, text=cpud] at ({14.6+(\k+0.5)*\cw},6.9) {\b};}
\node[note, font=\scriptsize, anchor=north west, text width=12.0cm] at (13.2,6.5) {%
  \textbf{S} START: SDA falls while SCL is high. Then the 7-bit address (0x1D = 0011101), R/W = 0 (write).
  \textbf{ACK}: the addressed device pulls SDA low. Data bytes follow, each with an ACK. \textbf{P} STOP: SDA rises while SCL is high.\\[3pt]
  Slow (100 kHz, 400 kHz, 1 MHz) but only 2 wires for many devices. Used for sensors, power chips (PMIC), cameras' control.
  The DE10-Lite's accelerometer can be read over I2C or SPI.};

% =====================================================================
% (F) GPIO, timer, watchdog
% =====================================================================
\panel{25.7}{0.7}{38.3}{14.7}{GPIO, timer and watchdog}
\node[anchor=north west, font=\scriptsize\bfseries] at (26.0,13.95) {GPIO: one bit per pin in each register};
\foreach \nm/\tx [count=\k from 0] in {DIR/{1 = output}, OUT/{value to drive}, IN/{after a 2-FF synchronizer}, IRQ\_EN/{interrupt on edge?}, IRQ\_ST/{edge seen (sticky)}}{
  \node[font=\tiny\ttfamily, anchor=west] at (26.0,{13.25-\k*0.5}) {\nm};
  \ffrow{27.4}{13.13-\k*0.5}{8}{0.24}{ff=prd}
  \node[lab, anchor=west] at (29.4,{13.25-\k*0.5}) {\tx};}
\node[note, font=\scriptsize, anchor=north west, text width=12.0cm] at (26.0,10.7) {%
  Edge detect = IN AND NOT IN\_previous: one flip-flop + one gate per pin. Your KEY buttons go through \rn{sync\_2ff} the same way.};
\node[anchor=north west, font=\scriptsize\bfseries] at (26.0,9.75) {timer / PWM};
\node[comb, minimum width=2.0cm, minimum height=0.7cm, font=\tiny] (pre) at (27.2,8.75) {prescaler $\div$k};
\node[draw=prd, fill=prf, thick, minimum width=2.4cm, minimum height=0.7cm, font=\tiny] (cnt) at (30.0,8.75) {32-bit counter (FF)};
\node[comb, minimum width=2.6cm, minimum height=0.7cm, font=\tiny] (cmp) at (33.1,8.75) {counter $<$ compare?};
\draw[dat] (pre) -- (cnt); \draw[dat] (cnt) -- (cmp);
\draw[dat] (cmp.east) -- (35.6,8.75); \node[lab, text=ink, anchor=west] at (35.6,8.75) {PWM pin / IRQ};
\renewcommand{\cw}{0.4}
\wname{27.4}{7.25}{PWM}  \wsig{27.4}{7.25}{1,1,0,0,0,0,0,0,1,1,0,0,0,0,0,0,1,1,0,0,0,0,0,0}
\node[note, font=\scriptsize, anchor=north west, text width=12.0cm] at (26.0,6.7) {%
  compare = 25\,\% of the period $\Rightarrow$ 25\,\% duty cycle (LED brightness, motor speed). A match can also raise an IRQ:
  that is how the OS gets its regular ``tick''. Your \rn{tick\_gen.sv} (2 Hz heartbeat LED) is a timer like this.};
\node[anchor=north west, font=\scriptsize\bfseries] at (26.0,5.1) {watchdog};
\node[note, font=\scriptsize, anchor=north west, text width=12.0cm] at (26.0,4.6) {%
  A counter that counts down all the time. Healthy software writes a ``kick'' register now and then, which reloads it.
  If the software hangs, the counter reaches 0 and the watchdog \textbf{resets the whole chip}. Cheap insurance:
  one counter, one comparator, one reset wire. Required in cars and industrial products.};

\figfoot{Addresses, register layouts and the I2C device addresses are examples. UART numbers match your design: 50 MHz, 115,200 baud, 434 clocks per bit.}
\end{tikzpicture}
\end{document}
